\documentclass{article}
\usepackage{graphicx}
\usepackage{amsmath}

\begin{document}

\begin{center}
\Large \textbf{avr-gcc Question}

\vspace{0.3cm}

Harshita N Kumar \\
COMETFWC052
\end{center}

\section*{Question}

For the output $F$ to be $1$ in the logic circuit shown, the input combination should be

\begin{itemize}
\item a) $A=1, B=1, C=0$
\item b) $A=1, B=0, C=0$
\item c) $A=0, B=1, C=0$
\item d) $A=0, B=0, C=1$
\end{itemize}

\section*{Solution}

By evaluating the output $F$ for each input combination using the logic function:

\begin{center}
\begin{tabular}{|c|c|c|c|}
\hline
A & B & C & F \\
\hline
1 & 1 & 0 & 0 \\
1 & 0 & 0 & 1 \\
0 & 1 & 0 & 0 \\
0 & 0 & 1 & 0 \\
\hline
\end{tabular}
\end{center}

Thus the output $F=1$ only for

\[
A=1,\ B=0,\ C=0
\]

Hence the correct option is

\begin{center}
\textbf{b) $A=1, B=0, C=0$}
\end{center}

\section*{Graph Output}

\begin{center}
\includegraphics[width=12cm]{graph.jpg}
\end{center}

\end{document}
